\documentclass[../main.tex]{subfiles}\begin{document}
\section{Basics}
Denne side indeholder alle vores basale indformationer omkring Echo-projektet.

\subsection{Ugentlig skema - DONE}

\begin{tabularx}{\textwidth}{
    >{\bfseries}p{2.5cm}
    p{2cm}
    X
}
\hline
Dag & Tid & Plan \\ 
\hline

Mandag 
& 08:00--16:00 
& Studiejob -- hele dagen \\

Tirsdag 
& 08:00--16:00 
& Uni-dag -- individuelle fag og gruppefag \\

Onsdag 
& 08:00--16:00 
& Studiejob -- hele dagen \\

Torsdag 
& 08:00--16:00 
& Bachelor-dag + møde med Torben \\

Fredag 
& 08:00--14:00 
& Padel + møde mellem boys + bachelor-dag \\

Lørdag 
& -- 
& Individuel dag \\

Søndag 
& -- 
& Individuel dag \\

\hline
\end{tabularx}

\subsection{Checklist - DONE}
\begin{tabularx}{\textwidth}{
    >{\bfseries}p{4cm}
    X
}
\hline
Område & Beskrivelse \\
\hline

Agile arbejdsværktøjer
&
GitHub Projects anvendes som det primære agile arbejdsværktøj. Dette inkluderer:
\begin{itemize}
    \item Kanban-board til opgavestyring
    \item Repository-dokumentation
\end{itemize}
\\

Tech-stack / risikoanalyse
&
Følgende teknologier overvejes og vurderes:
\begin{itemize}
    \item \textbf{Backend:} .NET/C\#, Golang, TypeScript/Node.js
    \item \textbf{Frontend -- Web:} React, Angular, Vue
    \item \textbf{Frontend -- Mobile:} Flutter (overvej her dart som sprog), .NET MAUI, SwiftUI + Kotlin
    \item \textbf{Database:} PostgreSQL, MongoDB
\end{itemize}
\\

Regler for produktudvikling
&
Der fastlægges fælles regler og standarder for udviklingsprocessen:
\begin{itemize}
    \item Gruppekontrakt
    \item Pull Requests (PRs)
    \item Projektopsætning
    \item Type checking
    \item Branch rules
    \item Genbrugelige komponenter
    \item Typografi
    \item Prettier / automatisk formattering
\end{itemize}
\\

\hline
\end{tabularx}

% Gantt-diagram
\begin{figure}[H]
    \centering
    \includegraphics[width=1\linewidth]{IMG/gantt_diagram.png}
    \caption{Gantt diagram for projektet}
    \label{fig:gantt}
\end{figure}

% What are we trying to solve???
\section{Problemformulering}
Hvordan kan en digital platform understøtte og motivere lokale borgere til at hjælpe hinanden gennem både betalte hverdagsopgaver og frivillige aktiviteter, samtidig med at platformen skaber tillid mellem aktørerne og muliggør verificering og belønning af sociale bidrag?

\subsection{Underspørgsmål}
1. Hvordan kan en fælles platform understøtte forskellige typer af hjælp, herunder både betalte peer-to-peer-opgaver og frivillige aktiviteter?\\

2. Hvordan kan platformen skabe tillid mellem brugere, organisationer og virksomheder, når aktørerne ikke nødvendigvis kender hinanden på forhånd?\\

3. Hvordan kan udførelsen af frivillige aktiviteter verificeres på en måde, der reducerer risikoen for misbrug af platformens belønningssystem?\\

4. Hvordan kan et pointbaseret belønningssystem designes, så brugere belønnes for sociale bidrag, samtidig med at transaktioner og indløsninger håndteres konsistent og sporbar?\\

5. Hvordan kan platformen designes, så forskellige aktørtyper og fremtidige typer af aktiviteter og belønninger kan tilføjes uden væsentlige ændringer af systemets kerne?\\

\subsection{Produktet + Aktører}

% Diagrams for clarity
\begin{figure}[H]
    \centering
    \includegraphics[width=1\linewidth]{IMG/actor.png}
    \caption{Aktører i produktet}
    \label{fig:actors}
\end{figure}


\section{Scope}
https://chatgpt.com/share/6a833c32-d56c-83ed-bdc5-0ba9b0e0081d



% Roles (frontend, backend, etc.) + tasks before semester start

\end{document}


% intro + rigt billede

% Kravspec (delt i epics, US og lign) (Epic aktor-kontekst diagram w. color-coding + tables and explanations)
% Functional and non-functional (FURPS+ in a table, focus on security and such)

% Metode og process (C4 og N+1 view) - skitser for hvordan det er brugt (GitHub projects)

% Analyse - Teknisk (FE, BE, DB osv) + Risiko + Brugerflade (UI/UX) and mock-ups

% Arkitektur (C4 + Domæne + applikations, osv) N+1
